\chapter*{РЕФЕРАТ}
\addcontentsline{toc}{chapter}{РЕФЕРАТ}
\thispagestyle{fancy}

% TODO: заповнити після завершення роботи

Робота містить \underline{\hspace{1cm}} сторінок, \underline{\hspace{1cm}} таблиць,
\underline{\hspace{1cm}} рисунків, \underline{\hspace{1cm}} джерел, \underline{\hspace{1cm}} додатків.

\textbf{Ключові слова:} VPN, WireGuard, двофакторна автентифікація, TOTP,
мережева безпека, захист даних, тунелювання трафіку.

\textbf{Мета роботи} --- розробка VPN-застосунку з двофакторною автентифікацією
на основі протоколу WireGuard, що поєднує захищене тунелювання мережевого
трафіку з надійним механізмом автентифікації користувача.

\textbf{Об'єкт дослідження} --- процес автентифікації користувачів у VPN-мережах.

\textbf{Предмет дослідження} --- методи та засоби інтеграції двофакторної
автентифікації з протоколом WireGuard.

\textbf{Методи дослідження}: аналіз та порівняння існуючих VPN-протоколів
і методів автентифікації; проєктування програмного забезпечення; тестування.

\textbf{Результати роботи}: розроблено програмний комплекс, що реалізує
VPN-з'єднання на основі WireGuard з обов'язковою двофакторною автентифікацією
за стандартом TOTP (RFC~6238).

\textbf{Наукова новизна:} удосконалено підхід до інтеграції двофакторної
автентифікації з протоколом WireGuard на основі моделі керування життєвим
циклом пірів, без модифікації протоколу та без використання попередньо
узгоджених ключів.

\newpage
